% Shared preamble for the three Weilight Harbor documents.
% Compile each main file with pdfLaTeX (Overleaf default).

\documentclass[11pt,a4paper]{article}

% --- Page layout
\usepackage[a4paper,margin=2.4cm]{geometry}
\usepackage{setspace}
\setstretch{1.10}
\setlength{\parskip}{4pt}
\setlength{\parindent}{0pt}

% --- Fonts and basics
\usepackage[utf8]{inputenc}
\usepackage[T1]{fontenc}
\usepackage{microtype}
\usepackage{lmodern}
\usepackage{textcomp}
% Map Unicode em/en-dashes and the simple bullet to LaTeX equivalents so they
% render in the default Computer Modern font without falling back to question
% marks.
\DeclareUnicodeCharacter{2014}{---}
\DeclareUnicodeCharacter{2013}{--}
\DeclareUnicodeCharacter{2026}{\ldots}
\DeclareUnicodeCharacter{2018}{`}
\DeclareUnicodeCharacter{2019}{'}
\DeclareUnicodeCharacter{201C}{``}
\DeclareUnicodeCharacter{201D}{''}
\DeclareUnicodeCharacter{2022}{\textbullet}

% --- Common packages
\usepackage{graphicx}
\usepackage{tabularx}
\usepackage{booktabs}
\usepackage{array}
\usepackage{longtable}
\usepackage{enumitem}
\setlist{nosep,leftmargin=*}
\usepackage{xcolor}
\usepackage{hyperref}
\hypersetup{
  colorlinks=true,
  linkcolor=black,
  citecolor=black,
  urlcolor=blue!55!black,
  pdfborder={0 0 0}
}
\usepackage{titlesec}
\titleformat{\section}{\large\bfseries}{\thesection}{0.6em}{}
\titleformat{\subsection}{\normalsize\bfseries}{\thesubsection}{0.5em}{}
\titleformat{\subsubsection}{\normalsize\itshape}{\thesubsubsection}{0.5em}{}
\titlespacing*{\section}{0pt}{1.6ex}{0.8ex}
\titlespacing*{\subsection}{0pt}{1.2ex}{0.5ex}
\titlespacing*{\subsubsection}{0pt}{1.0ex}{0.3ex}

% Inline code
\usepackage{listings}
\lstset{
  basicstyle=\ttfamily\footnotesize,
  breaklines=true,
  columns=fullflexible,
  showstringspaces=false,
  frame=single,
  framerule=0.3pt,
  rulecolor=\color{black!30},
  xleftmargin=0pt,
  xrightmargin=0pt,
  aboveskip=4pt,
  belowskip=4pt
}

% Tighter table column types
\newcolumntype{L}[1]{>{\raggedright\arraybackslash}p{#1}}
\renewcommand{\arraystretch}{1.15}

% Header / footer
\usepackage{fancyhdr}
\pagestyle{fancy}
\fancyhf{}
\fancyfoot[C]{\thepage}
\fancyhead[L]{\small Weilight Harbor — Group 14}
\fancyhead[R]{\small COMP3030J 2025/26}
\renewcommand{\headrulewidth}{0.3pt}


\title{\vspace{-1.5cm}\textbf{Weilight Harbor — User Document}\\
       \large COMP3030J Web Development, UCD, 2025/26 — Group 14 (Draft, Week 11)}
\author{}
\date{}

\begin{document}
\maketitle
\vspace{-2.0cm}
\begin{center}
  \small
  Live deployment: \url{http://csi6220-2-vm3.ucd.ie}
\end{center}
\vspace{0.4cm}

\noindent
This document is written for the three end-user roles supported by
Weilight Harbor. Screenshots are referenced inline; the placeholder
captions will be replaced with actual screenshots before final
submission.

% =====================================================================
\section{Getting started}

\subsection{What is Weilight Harbor?}
Weilight Harbor (\textit{Weiguang Gangwan}) is a charity and mutual-aid platform
for the families and caregivers of seriously ill patients. It combines a
moderated community, an after-dark anonymous space, a private journal, a
verified crowdfunding module, and an interactive map for arranging
temporary respite help. The platform deliberately puts emotional
support \emph{ahead} of financial transaction.

\subsection{Browser and device support}
\begin{itemize}
  \item Recommended desktop browsers: latest Chrome, Firefox, Safari, or
        Edge.
  \item Mobile: iOS Safari $\geq 16$, Android Chrome $\geq 120$.
  \item A network connection is required to reach the live VM.
\end{itemize}

\subsection{Visiting the site}
Open \url{http://csi6220-2-vm3.ucd.ie} in a browser. The landing page
introduces the platform; the navigation bar at the top lets you move
between modules, switch language (\textbf{ZH}/\textbf{EN}), and log in.

\subsection{Test accounts (assignment evaluation only)}
\begin{tabularx}{\linewidth}{@{} l l L{4cm} X @{}}
  \toprule
  \textbf{Username} & \textbf{Password} & \textbf{Role} & \textbf{Use this for} \\
  \midrule
  \texttt{admin}    & \texttt{admin123} & Administrator       & Admin
                                                                panel,
                                                                review
                                                                queues \\
  \texttt{limei}    & \texttt{test123}  & Certified Family    & Verified
                                                                caregiver
                                                                journey \\
  \texttt{zhangwei} & \texttt{test123}  & Regular User        & Donor /
                                                                visitor
                                                                journey \\
  \bottomrule
\end{tabularx}

\noindent
A complete account list is in Appendix~A.

\subsection{Switching language}
Use the \textbf{ZH}/\textbf{EN} toggle on the left of the navigation
bar. The choice is remembered for the rest of your session and across
visits on the same browser.

% =====================================================================
\section{Regular User}

A \emph{regular user} is anyone who has registered without applying for
caregiver certification. Regular users can browse, donate, post in the
community, comment, react, and use the night confession panel.

\subsection{Registering and logging in}
\begin{enumerate}
  \item Click \textbf{Join the Harbor} in the top-right of any page.
  \item Fill in username, email or phone, password ($\geq$~6 characters),
        and confirm password.
  \item Submit. You are redirected to the login page; sign in with your
        credentials.
  \item Forgot your password? Click \textbf{Forgot password?} on the
        login page (UI only in the current build).
\end{enumerate}

\textit{[Screenshot: Register page]}

\subsection{Browsing the community}
The community page (\texttt{/community/}) is split: the left 70\,\% is
the post feed; the right 30\,\% is the \textbf{Night Confessions} panel.
Filter posts by the category pills near the top (e.g.\ \emph{Emotional
Support}, \emph{Daily Life}, \emph{Medical Info}, \emph{Good News},
\emph{Resources}). Sort by \textbf{Latest}, \textbf{Most Empathised}, or
\textbf{Hot}. Click a post card to open the detail page.

\textit{[Screenshot: Community feed with right-side Night Confessions
panel]}

\subsection{Reacting and commenting on a post}
\begin{itemize}
  \item The empathy button (heart icon) is labelled \emph{"I feel the
        same"} and toggles your reaction. Each user can react once per
        post.
  \item Comments appear chronologically; click \textbf{Reply} under a
        comment to thread.
  \item You can mark your comment \textbf{Anonymous} via the checkbox
        above the comment box.
\end{itemize}

\subsection{Writing a post}
\begin{enumerate}
  \item Click \textbf{New Post} in the community navigation.
  \item Choose at least one category, optionally add a title
        (max.~50 characters), write your content, attach up to nine
        images.
  \item Toggle \textbf{Post anonymously} if you prefer; your username
        and avatar will be hidden.
  \item Click \textbf{Publish}. If your text contains words flagged as
        sensitive, you will be warned; if the text contains crisis
        keywords, you will be redirected to a hotline modal before
        submission.
\end{enumerate}

\subsection{Night Confessions}
\begin{itemize}
  \item The panel is \textbf{active 21:00–06:00} server local time
        (UTC on the VM). Outside that window the panel is dimmed with a
        countdown.
  \item During active hours you can write a confession (anonymous,
        capped length) or send a \textbf{hug} reaction to an existing
        one.
  \item If your text mentions self-harm or crisis keywords, a crisis
        hotline modal appears \emph{before} submission.
\end{itemize}

\textit{[Screenshot: Night Confessions modal at night]}

\subsection{Donating to a crowdfunding campaign}
\begin{enumerate}
  \item Open \texttt{/crowdfunding/}.
  \item Filter or search for a campaign.
  \item Click a campaign card to open the detail page.
  \item The right-hand sticky panel shows progress and the
        \textbf{Donate} button. Choose a preset amount or enter a custom
        value.
  \item Optionally write a public message and tick \textbf{Anonymous} to
        hide your name from the donor list.
  \item Confirm. The progress bar animates and a thank-you toast
        appears. \emph{Donations are simulated for the assignment.}
\end{enumerate}

\subsection{Personal centre}
\begin{itemize}
  \item \textbf{Profile} (\texttt{/user/profile}) — overview of your
        activity.
  \item \textbf{Edit profile} (\texttt{/user/edit}) — name, phone,
        email, avatar, bio.
  \item \textbf{History} (\texttt{/user/history}) — paginated tabs for
        your Posts, Campaigns, Donations and Respite Requests.
  \item \textbf{Settings} (\texttt{/user/settings}) — change password.
  \item \textbf{Apply for certification} (\texttt{/user/certification})
        — see Section~3.1.
\end{itemize}

\subsection{Logging out}
Click your avatar in the top-right $\to$ \textbf{Logout}. Logout is
issued as a POST so it cannot be triggered from a malicious link.

% =====================================================================
\section{Certified Family / Volunteer}

A certified user is a regular user whose caregiver/volunteer status has
been approved by an administrator. Certified users keep all
regular-user capabilities and additionally gain access to crowdfunding
creation, the journal, and the create/accept actions on the respite
map.

\subsection{Applying for certification}
\begin{enumerate}
  \item Open the avatar menu $\to$ \textbf{Apply Certification}
        (\texttt{/user/certification}).
  \item Choose \textbf{Family} or \textbf{Volunteer}.
  \item Fill in real name, ID number, relationship to the patient, and
        patient details (name, illness, hospital, diagnosis date).
  \item Upload a primary proof document (PDF / JPG / PNG). Optionally
        upload up to three additional supporting documents.
  \item Submit. Your application enters a queue with status
        \textbf{Pending}; an admin will approve or reject within
        24–48 hours.
  \item Track status at \texttt{/user/certification/status}. If
        rejected, the rejection reason is shown; you can re-apply.
\end{enumerate}

\textit{[Screenshot: Certification status page with timeline]}

\subsection{Starting a crowdfunding campaign}
\begin{enumerate}
  \item From the crowdfunding page, click \textbf{Start a Campaign}.
  \item If you are not yet certified, the page redirects to the
        certification application.
  \item Provide title, story, target amount, deadline, category,
        cover image, patient photo and a payment QR placeholder.
  \item Click \textbf{Submit for Review}. Status becomes \textbf{Pending};
        once approved by an admin, your campaign appears in the public
        listing.
\end{enumerate}

\subsection{Receiving donations}
\begin{itemize}
  \item The campaign detail page shows real-time progress, donor count,
        and a donor list (anonymous donations render as \emph{"Anonymous
        Caregiver"}).
  \item When \texttt{current\_amount} $\geq$ \texttt{funding\_goal}, a
        \textbf{Fully Funded} badge is shown and the donate button is
        greyed out.
\end{itemize}

\subsection{Life Journal}
\begin{itemize}
  \item Open \texttt{/journal/} for your private timeline. Only you can
        see your entries.
  \item \textbf{New entry} — pick a date, a mood emoji, write your
        text, attach up to nine photos.
  \item \textbf{Edit / Delete} — available from any entry detail page
        (soft-delete preserves the row).
  \item \textbf{Filters} — choose a year and month to focus the
        timeline.
  \item \textbf{Annual Film} — once you have at least ten entries in a
        year, an \emph{Annual Film} button unlocks; it plays a
        full-screen Ken-Burns slideshow of up to twelve highlights.
\end{itemize}

\textit{[Screenshot: Journal timeline + entry composer]}

\subsection{Respite Station — posting and accepting}
\begin{enumerate}
  \item Open \texttt{/respite/}. The map centres on Dublin by default;
        allow location access for a closer view.
  \item Click \textbf{I need a break} to post a service request, or
        \textbf{I can lend equipment} to post an equipment offer.
  \item Fill the form: title, time window, address (auto-filled if
        location is allowed), special notes, contact preference.
  \item Your pin appears on the map (orange = service, green =
        equipment) with a pulsing animation while pending.
  \item Other certified users can click your pin and \textbf{Accept}.
        You cannot accept your own request.
  \item After the activity is over, the original requester clicks
        \textbf{Confirm Completed} to close the request.
\end{enumerate}

\subsection{Profile visibility}
Certified users display a verification badge next to their name across
the platform: in posts, comments, donations, and respite pins.

% =====================================================================
\section{Administrator}

Administrators see an additional \textbf{Admin} entry in the user-menu
dropdown that opens \texttt{/admin/}. The admin panel is a Layui-based
interface inherited from Pear-Admin-Flask, extended with three custom
review pages.

\subsection{Logging in as an administrator}
Use the seeded \texttt{admin}/\texttt{admin123} account on the standard
\texttt{/auth/login} page. After signing in, your avatar dropdown
contains an \textbf{Admin Panel} item, which opens \texttt{/admin/} in a
new tab.

\subsection{Dashboard}
The admin dashboard shows summary cards (pending counts) and a left-hand
sidebar with sections:
\begin{itemize}
  \item \textbf{Certification Reviews} — \texttt{/admin/certification/}
  \item \textbf{Campaign Reviews} — \texttt{/admin/campaign-review/}
  \item \textbf{Community Moderation} — \texttt{/admin/community/}
  \item \textbf{User Management} — inherited from the base panel
  \item \textbf{Roles \& Powers} — inherited from the base panel
  \item \textbf{Operation Log} — inherited audit log
\end{itemize}

\subsection{Reviewing certification applications}
\begin{enumerate}
  \item Open \textbf{Certification Reviews}.
  \item The table lists applicants oldest first; filter by status if
        needed.
  \item Click \textbf{Review} on a row to see all uploaded documents
        (clickable to enlarge), the applicant profile, and any prior
        applications.
  \item Pick one of:
        \begin{itemize}
          \item \textbf{Approve} — confirms the certification; the
                applicant is granted certified status immediately.
          \item \textbf{Reject} — requires a reason; shown to the
                applicant on their status page.
          \item \textbf{Request More Info} — sends a note back; status
                stays \emph{Pending}.
        \end{itemize}
  \item Every action is logged to the operation log table.
\end{enumerate}

\textit{[Screenshot: Certification review modal with documents]}

\subsection{Reviewing crowdfunding campaigns}
Open \textbf{Campaign Reviews}. Click a campaign to open the review
pane: cover image, story, payment QR, target amount and any patient
documents. \textbf{Approve} to publish the campaign;
\textbf{Reject} with a reason. The campaign creator sees the decision on
their personal centre.

\subsection{Moderating community posts}
\begin{enumerate}
  \item Open \textbf{Community Moderation}.
  \item Filter by status (auto, hidden, pending) or search by author or
        keyword.
  \item For each post you can:
        \begin{itemize}
          \item \textbf{Approve} — leave it visible.
          \item \textbf{Hide} — keep the row in the database but remove
                it from the public feed; the author sees a \emph{"hidden
                by moderator"} notice.
          \item \textbf{Delete} — soft-delete (sets \texttt{delete\_at});
                reversible only by direct database operation.
        \end{itemize}
\end{enumerate}

\subsection{User and role management}
The base Pear-Admin-Flask \textbf{User Management} and \textbf{Role /
Power} pages remain available for managing accounts, freezing users, or
editing the RBAC matrix. New roles or permissions follow the existing
dictionary structure.

\subsection{Audit log}
Every action wrapped in the \texttt{@authorize("…", log=True)}
decorator writes a row to the operation log, including the operator,
the resource code, and a timestamp.

% =====================================================================
\section{Frequently Asked Questions}

\textbf{Q1. The Night Confessions panel is dimmed — is the feature
broken?}\\
No. It is intentionally only active between 21:00 and 06:00 (VM local
time, UTC).

\textbf{Q2. Why does the \emph{Donate} button do nothing for my
campaign?}\\
The campaign is probably already fully funded, or your account is not
signed in. The button is greyed out in both cases.

\textbf{Q3. I uploaded a document but the page says "valid proof
document required".}\\
The file may exceed the size limit or have a non-permitted extension.
Allowed types are PDF, JPG, JPEG, PNG (and the equivalents declared in
\texttt{ALLOWED\_DOC\_EXTS}).

\textbf{Q4. Why can my friend not see my journal entries?}\\
The journal is private by design. All journal queries filter by
\texttt{user\_id == current\_user.id}; no one but you, including
administrators, can read your entries through the public site.

\textbf{Q5. I cannot start a campaign — the page tells me to "apply for
certification".}\\
Crowdfunding requires verified caregiver/volunteer status. Submit a
certification application (Section~3.1) and wait for admin approval
(24–48 hours).

\textbf{Q6. The map does not centre on my location.}\\
The browser blocked the location prompt or no permission was granted.
The map falls back to Dublin city centre. You can pan/zoom manually.

\textbf{Q7. Where do I report a bug?}\\
For the assignment evaluation, please email the group contact listed on
the cover of the system document, or open an issue at
\url{https://github.com/licctvcctv/weilight-harbor/issues}.

% =====================================================================
\appendix
\section*{Appendix A — Test accounts (full list)}
\addcontentsline{toc}{section}{Appendix A — Test accounts}
{\small
\begin{tabularx}{\linewidth}{@{} l l L{4cm} X @{}}
  \toprule
  \textbf{Username} & \textbf{Password} & \textbf{Role} & \textbf{Notes} \\
  \midrule
  \texttt{admin}     & \texttt{admin123} & Administrator       & Full admin access \\
  \texttt{limei}     & \texttt{test123}  & Certified Family    & ALS caregiver, has campaigns \\
  \texttt{zhangwei}  & \texttt{test123}  & Regular User        & Donor profile, tech background \\
  \texttt{wangfang}  & \texttt{test123}  & Certified Volunteer & Retired nurse \\
  \texttt{chenyun}   & \texttt{test123}  & Certified Family    & Alzheimer's caregiver \\
  \texttt{liuyang}   & \texttt{test123}  & Regular User        & Has a pending certification \\
  \texttt{sunli}     & \texttt{test123}  & Certified Family    & Pediatric cancer caregiver \\
  \texttt{zhaojing}  & \texttt{test123}  & Certified Volunteer & Social worker \\
  \bottomrule
\end{tabularx}}

\section*{Appendix B — Glossary}
\addcontentsline{toc}{section}{Appendix B — Glossary}
\begin{itemize}
  \item \textbf{Certified user} — a regular user whose caregiver or
        volunteer status has been approved by an administrator.
  \item \textbf{Empathy reaction} — the like equivalent on this
        platform, labelled \emph{"I feel the same"} (wo ye shi).
  \item \textbf{Night Confessions} — the time-gated, anonymous
        after-dark panel.
  \item \textbf{Annual Film} — a full-screen slideshow synthesised from
        a user's journal entries.
  \item \textbf{Respite} — a temporary break from caregiving
        responsibilities.
\end{itemize}

\end{document}
